%%%%%%%%%%%%%%%%%%%%%%%%%%%%%%%%%%%%%%%%%
% a0poster Portrait Poster
% LaTeX Template
% Version 1.0 (22/06/13)
%
% The a0poster class was created by:
% Gerlinde Kettl and Matthias Weiser (tex@kettl.de)
% 
% This template has been downloaded from:
% http://www.LaTeXTemplates.com
%
% License:
% CC BY-NC-SA 3.0 (http://creativecommons.org/licenses/by-nc-sa/3.0/)
%
%%%%%%%%%%%%%%%%%%%%%%%%%%%%%%%%%%%%%%%%%

%----------------------------------------------------------------------------------------
%	PACKAGES AND OTHER DOCUMENT CONFIGURATIONS
%----------------------------------------------------------------------------------------

\documentclass[a0,portrait]{a0poster}

\usepackage{multicol} % This is so we can have multiple columns of text side-by-side
\columnsep=100pt % This is the amount of white space between the columns in the poster
\columnseprule=3pt % This is the thickness of the black line between the columns in the poster

\usepackage[svgnames]{xcolor} % Specify colors by their 'svgnames', for a full list of all colors available see here: http://www.latextemplates.com/svgnames-colors

\usepackage{times} % Use the times font
%\usepackage{palatino} % Uncomment to use the Palatino font

\usepackage{graphicx} % Required for including images
\graphicspath{{figures/}} % Location of the graphics files
\usepackage{booktabs} % Top and bottom rules for table
\usepackage[font=small,labelfont=bf]{caption} % Required for specifying captions to tables and figures
\usepackage{amsfonts, amsmath, amsthm, amssymb} % For math fonts, symbols and environments
\usepackage{wrapfig} % Allows wrapping text around tables and figures

\renewcommand\refname{Referenzen}

\begin{document}

%----------------------------------------------------------------------------------------
%	POSTER HEADER 
%----------------------------------------------------------------------------------------

% The header is divided into two boxes:
% The first is 75% wide and houses the title, subtitle, names, university/organization and contact information
% The second is 25% wide and houses a logo for your university/organization or a photo of you
% The widths of these boxes can be easily edited to accommodate your content as you see fit

\begin{minipage}[b]{0.75\linewidth}
\veryHuge \color{NavyBlue} \textbf{Die MultiLA Softwareplattform} \color{Black}\\ % Title
\Huge\textit{Authoring Software und Learning Analytics für interaktive Lernanwendungen im Bereich Statistik und Data Science}\\[2cm] % Subtitle
\huge Markus Konrad MSc.\textsuperscript{1,2},
Prof. Dr. Maria Osipenko\textsuperscript{2},
Prof. Dr. Martin Spott\textsuperscript{1},
Prof.~Dr.~Andre Beinrucker\textsuperscript{1}\\[0.5cm] % Author(s)
\large \textsuperscript{1}Hochschule für Technik und Wirtschaft Berlin, \textsuperscript{2}Hochschule für Wirtschaft und Recht Berlin\\[0.4cm] % University/organization
%\Large \texttt{john@LaTeXTemplates.com} --- 1 (000) 111 1111\\
\end{minipage}
%
\begin{minipage}[b]{0.25\linewidth}
% TODO: QR Code hier
% TODO: HTW, HWR, IFAF Logos ganz unten
\includegraphics[width=20cm]{logo.png}\\
\end{minipage}

\vspace{1cm} % A bit of extra whitespace between the header and poster content

%----------------------------------------------------------------------------------------

\begin{multicols}{2} % This is how many columns your poster will be broken into, a portrait poster is generally split into 2 columns

%----------------------------------------------------------------------------------------
%	ABSTRACT
%----------------------------------------------------------------------------------------

\color{Navy} % Navy color for the abstract

\begin{abstract}

TODO

\end{abstract}

%----------------------------------------------------------------------------------------
%	INTRODUCTION
%----------------------------------------------------------------------------------------

\color{DarkSlateGray} % DarkSlateGray color for the rest of the content

\section*{Einleitung}

Mittels der \textit{MultiLA Softwareplattform} ist es möglich, web-basierte, interaktive Lernanwendungen zu entwickeln, diese den Studierenden zur Verfügung zu stellen und auf datenschutzkonformeweise Interaktionsdaten der Studierenden mit den Lernanwendungen zu sammeln sowie Experimente durchzuführen. Damit werden drei Ziele erreicht:

\begin{itemize}
    \item Verbesserung der Lehre durch innovative, interaktive Lernanwendungen,
    \item Unterstützung der Forschung im Bereich Learning Analytics,
    \item Verbesserung der Lernanwendungen durch den aus den Daten gewonnenen Erkenntnissen.
\end{itemize}

Die Softwareplattform wurde im Rahmen des IFAF MultiLA-Projekts an der HTW Berlin und der HWR Berlin entwickelt und steht komplett als Open Source Softwarepaket zur Verfügung.

\subsection*{Features}

\begin{itemize}
    \item Erstellung interaktiver Lernanwendungen mittels \textit{RMarkdown} oder als \textit{Shiny}-Anwendungen
    \item Hoch granulare, konfigurierbare und anonyme Nachverfolgung von Benutzerinteraktionen mit den Lernanwendungen: Mausbewegungen, Klicks, Einreichen von Übungen usw.
    \item Unterstützung von A/B-Testexperimenten und integrierten Umfragen
    \item Konfigurierbare Lernanwendungen: Basisanwendungen einmal schreiben, Varianten über Konfigurationen erstellen
    \item Dynamische Zusammenfassungsleiste für Lernanwendungen
    \item Webbasierte Administrationsoberfläche zum Veröffentlichen von Lernanwendungen, Einrichten von Varianten und Experimenten sowie Herunterladen gesammelter Daten
    \item Datenaufbereitungs- und Analyseskripte
    \item Datenschutzkonforme, selbst gehostete Lösung
\end{itemize}

\section*{Von der Lernanwendung zu Learning Analytics}

\subsection*{1. Lernanwendungen erstellen}

Der inhaltliche Rahmen liegt auf dem Bereich Statistik, Mathematik und Data Science. In diesem Bereich ist die Nutzung der Programmiersprache \textit{R} sowohl unter Lehrenden als auch in den Curricula sehr ausgeprägt, weshalb zur Erstellung der Lernanwendungen \textit{R-Markdown} in Kombination mit \textit{RStudio} zum Einsatz kommen.

\subsection*{2. Lernanwendungen bereitstellen}

TODO

\subsection*{3. Nutzungsdaten sammeln und auswerten}

TODO

\section*{Hintergrund}

\subsection*{Softwarearchitektur}

TODO

\subsection*{Erstellte Lernanwendungen}

TODO

\subsection*{Erste Auswertungen}

TODO

\section*{Fazit und Ausblick}


\begin{center}\vspace{1cm}
\includegraphics[width=0.8\linewidth]{placeholder}
\captionof{figure}{\color{Green} Figure caption}
\end{center}\vspace{1cm}

 %----------------------------------------------------------------------------------------
%	REFERENCES
%----------------------------------------------------------------------------------------

\nocite{*} % Print all references regardless of whether they were cited in the poster or not
\bibliographystyle{plain} % Plain referencing style
\bibliography{sample} % Use the example bibliography file sample.bib


\end{multicols}
\end{document}